\documentclass[UTF8,zihao=-4,fontset=none,a4paper]{ctexart}
\usepackage[left=2.5cm,right=2.5cm,top=2.5cm,bottom=2.4cm]{geometry}
\usepackage{fontspec}
\setCJKmainfont{SimSun}[AutoFakeBold=2.2]
\setCJKsansfont{SimSun}[AutoFakeBold=2.2]
\setCJKmonofont{SimSun}
\setmainfont{Times New Roman}
\setsansfont{Times New Roman}
\setmonofont{Times New Roman}
\usepackage{amsmath,amssymb,booktabs,tabularx,array,graphicx,float}
\usepackage{tikz}
\usetikzlibrary{arrows.meta,positioning,calc,fit}
\usepackage{xcolor,setspace,fancyhdr,caption,hyperref,enumitem,listings,tocloft,placeins}
\hypersetup{hidelinks}
\setstretch{1.32}
\setlength{\parindent}{2em}
\setlength{\parskip}{2pt}
\setcounter{tocdepth}{2}
\pagestyle{fancy}\fancyhf{}\fancyfoot[C]{\thepage}
\renewcommand{\headrulewidth}{0pt}\renewcommand{\footrulewidth}{0pt}
\ctexset{section={format=\zihao{3}\bfseries\color{black},beforeskip=12pt,afterskip=16pt},subsection={format=\zihao{4}\bfseries\color{black},beforeskip=10pt,afterskip=7pt}}
\captionsetup{font=normalsize,labelfont=bf,skip=8pt}
\setlist{nosep,leftmargin=2em}
\newcolumntype{Y}{>{\raggedright\arraybackslash}X}
\newcommand{\pagechapter}[1]{\section{#1}}
\newcommand{\diagram}[2]{\begin{figure}[htbp]\centering\input{figures/#1.tex}\caption{#2}\end{figure}}
\newcommand{\screen}[2]{\begin{figure}[htbp]\centering\includegraphics[width=\linewidth]{images/#1.png}\caption{#2}\end{figure}}
\lstset{basicstyle=\normalsize\ttfamily,keywordstyle=\color{black}\bfseries,commentstyle=\color{black},stringstyle=\color{black},breaklines=true,frame=single,rulecolor=\color{black!35},columns=fullflexible,keepspaces=true,showstringspaces=false}

\definecolor{celadon}{HTML}{2E6A67}
\definecolor{pale}{HTML}{D7E6DB}
\definecolor{tea}{HTML}{7E6554}
\definecolor{rice}{HTML}{E9D6BE}
\definecolor{fog}{HTML}{607E95}
\definecolor{silver}{HTML}{E5EBEF}
\definecolor{stone}{HTML}{6C7F4F}
\setlength{\textfloatsep}{12pt}
\setlength{\intextsep}{12pt}
\renewcommand{\topfraction}{.95}
\renewcommand{\bottomfraction}{.95}
\renewcommand{\textfraction}{.05}
\renewcommand{\floatpagefraction}{.85}
\setcounter{totalnumber}{5}
\renewcommand{\cftsecfont}{\bfseries}
\setlength{\cftbeforesecskip}{4pt}
\setlength{\cftbeforesubsecskip}{0pt}
\setlength{\cftsubsecindent}{1.5em}
\begin{document}\color{black}
\begin{titlepage}\centering
\vspace*{1.3cm}
{\zihao{3}\bfseries 第九届传智杯全国大学生数智创新与AI应用大赛\par}
\vspace{.6cm}{\zihao{4}AI创新应用挑战赛\par}
\vspace{2.5cm}
{\zihao{0}\bfseries 学脉 XueMai\par}
\vspace{.6cm}{\zihao{2}\bfseries 师生双端AI个性化智慧教育平台\par}
\vspace{1.5cm}{\zihao{2}作品技术文档\par}
\vfill
\begin{tabular}{rl}
作品名称：&\underline{\makebox[9cm][c]{学脉 XueMai}}\\[15pt]
作品类型：&\underline{\makebox[9cm][c]{师生双端AI个性化智慧教育平台}}\\[15pt]
参赛赛道：&\underline{\makebox[9cm][c]{AI创新应用挑战赛}}\\[15pt]
应用场景：&\underline{\makebox[9cm][c]{课程学习、课后巩固与教学协作}}\\[15pt]
系统版本：&\underline{\makebox[9cm][c]{0.5}}\\[15pt]
文档日期：&\underline{\makebox[9cm][c]{2026年11月}}
\end{tabular}
\vspace{1.6cm}
\end{titlepage}
{\fontsize{10.5pt}{12.6pt}\selectfont\setstretch{1}\tableofcontents}
\clearpage\section*{摘要}\addcontentsline{toc}{section}{摘要}
从课堂学习到课后巩固，学生遇到的困难往往分散在不同位置：一道题暴露出概念混淆，一段视频留下未解的问题，一次课堂回答反映了先修知识缺口。教师看到这些材料后，还需要判断应当帮助谁、安排什么内容，以及后续学习是否出现变化。学脉围绕这一过程，把课程知识结构、学生作答、AI辅导和教师任务组织在同一套教学平台中。

学生端以当前学习任务为入口，连接知识图谱、讲义、视频片段、交互实验和习题。系统根据题目难度、提示使用、部分得分、相近题相关性与时间间隔更新知识状态，同时呈现证据充分度。对于需要巩固的知识，平台沿先修关系组织学习路径，并把复习时间安排到学生的学习工作区。

教师端围绕课程组织、题库审核、主观题评阅和针对性任务展开。教师可以从知识点进入学生的具体作答，查看状态变化的计算依据，在发布任务时保存基线，再对照学生完成任务后的表现。诊断、策略和校验三个角色协作，AI负责课程资料解释、渐进提示和教学草稿，客观题评分与业务动作执行由后端控制。

系统采用React、TypeScript、FastAPI与SQLite实现，接入DeepSeek官方接口并提供本地示例模式。演示采用数据结构与Python两门课程，分别包含31个与27个知识点、28人与14人的合成教学班，以及讲义、多题型习题、课堂历史和六段原创教学视频。验证围绕算法性质、接口流程和师生浏览器操作开展，覆盖提示差异、重复作答、时间折减、间隔复习、辅导理解检查与教学任务回流。

学脉的主要特点是将“下一步学习内容”落实为能够直接进入的课程活动，并使学生状态和教师行动共享可追溯的学习证据。这为课程学习和日常教学提供了一个可运行、可解释且便于逐步完善的个性化教学方案。

\vspace{1em}\noindent\textbf{关键词：}个性化学习；知识图谱；学习证据；间隔复习；生成式人工智能


\section{应用场景与总体方案}

\subsection{应用场景与设计目标}

\subsubsection{课程学习中的信息断点}
数据结构与算法课程同时涉及概念、程序步骤和问题求解。学生能够背出“后进先出”，并不一定能正确追踪交错的入栈出栈操作；能够看懂一次排序动画，也不一定能说明下一步为何交换。课程平台如果只保存观看时长和题目得分，教师很难进一步判断困难发生在哪里，学生也容易在大量材料之间反复切换。

学脉把课程知识点作为连接材料与活动的共同对象。讲义、视频、题目和实验归属到具体知识点，作答则记录题目、答案、提示和时间。这样，学生从一道错题进入辅导时，可以继续讨论同一个知识点；教师查看困难知识时，也能够追溯到形成判断的具体作答。
\subsubsection{面向行动组织个性化学习}
平台的个性化输出是一个可执行的任务序列。推荐首先说明当前需要巩固的内容，再给出知识卡或视频片段、针对性练习、理解检查和后续复测入口。有先修缺口时，序列把前置知识放在前面，而不是直接增加目标题目的数量。

学生首页优先回答“现在学什么、为什么、如何继续”。教师首页则优先呈现课程操作、待复核作答和需要关注的知识点。两端通过同一套课程与学习记录协作，同时保持各自任务所需的信息密度。
\subsubsection{设计目标}
项目设置三个相互关联的目标。第一，使每个学习建议都能进入明确的课程对象，减少材料检索和页面切换的成本。第二，使用有明确定义的状态估计与复习规则，保留计算依据，使学习建议具有可解释性。第三，使教师的任务发布与学生的完成结果能够相互对应，便于跟踪后续学习表现。

平台以课程、章节、知识点和学习活动组织内容，可承载不同学段的学科学习、职业技能训练与课后巩固。教师根据教学目标配置知识依赖、题目与资源，学生沿相应课程开展学习。本文采用数据结构与Python课程展示具体过程，它们是便于观察概念、步骤与迁移差异的演示内容。个性化教学关注学习者的具体反馈，反馈质量与后续行动之间的联系也是教学研究的重要议题\cite{bloom,hattie,vanlehn}。


\subsection{角色分工与学习活动组织}

学生端围绕学习活动组织导航。学习概览显示下一项任务；目标与计划保存学生的个人安排；知识图谱展示概念之间的先修关系；练习、复习、视频和实验用于形成不同类型的学习经历。AI辅导位于这些活动之间，接收题目或视频片段的上下文，并提供相应教学动作。

教师端把课程建设和教学判断放在主要位置。教师维护章节、知识点、先修关系、讲义与视频，审核AI生成的题目和教案，发布面向班级或指定学生的任务。对于主观作答，教师依据量规确认分数与评语，再将结果写入学习记录。

\begin{table}[H]\centering\caption{两类角色的主要任务}\begin{tabularx}{\linewidth}{p{2cm}YY}\toprule 角色&主要输入&形成的输出\\\midrule 学生&课程材料、当前任务、辅导提示&答案、学习笔记、实验记录、个人安排\\教师&知识结构、学习证据、教学目标&课程内容、审核结果、评语、针对性任务\\AI助手&课程上下文、资料片段、用户问题&解释、提示、教学草稿、受支持动作选择\\\bottomrule\end{tabularx}\end{table}

同一个教学对象在不同角色中具有不同的操作方式。例如，学生查看知识点时关心掌握状态、先修与学习材料；教师查看同一节点时关心知识描述、资源组织和先修边的编辑。系统共享对象标识，避免两端形成彼此独立的材料副本。

课堂互动补充了课后学习之外的反馈。教师可以发起即时问题，学生提交自己的回答，教师端持续更新回答情况。课堂回答与课后练习各自保留语义，便于教师把课堂观察和后续作答结合起来使用。

个人计划既可由对应学习活动自动完成，也可由学生自报完成；知识状态根据实际作答证据更新。这种分工使“我完成了安排”和“我能独立解决问题”分别表达，给学生保留规划自主性，也使教师能够准确理解不同记录的作用。


\subsection{系统总体架构}

系统采用模块化单体架构。前端负责交互、课程对象定位和结果呈现；后端统一处理身份、课程访问范围、确定性评分、学习证据与AI调用。SQLite保存课程内容、学习事件和版本信息。对于本机体验和课程规模的演示，这一方案能够以较少的运行依赖贯通主要业务。
\diagram{architecture}{学脉总体架构与模块关系}
前端按学习和教学任务分成独立页面，通过路由按需加载。知识图谱布局、算法实验、富文本渲染和理解检查分别封装，减少页面之间的重复实现。主题样式由共享设计变量控制，学生和教师使用不同主色与页面结构。

后端将课程管理、练习评阅、资源、学生计划、教学分析和AI服务分开实现。学习证据引擎提供统一的状态查询与更新方法，首页、图谱、推荐和教师分析使用同一套计算规则。原始作答、证据变化和复习安排分别保存，以支持过程追踪。

AI适配层区分本地示例与官方服务。两种方式使用相同的会话和教学动作接口；在使用本地示例时，界面集中显示相应模式。实际AI解释与业务执行分离，模型提供的动作类型经校验后，才映射到课程内的具体对象。


\section{课程知识与渐进式图谱}

\subsection{课程知识结构与内容建设}

课程知识结构采用有向无环图。节点代表知识点，先修边从基础知识指向依赖它的后续内容。章节提供主题组织，知识点保存定义、难度和学习目标，资源与题目通过知识点关联。修改先修关系时，后端检查是否形成循环，使课程路径保持可解释的顺序。

\begin{table}[H]\centering\caption{数据结构与算法课程的六个主题}\begin{tabularx}{\linewidth}{p{2.5cm}Y}\toprule 主题&知识点\\\midrule 线性结构&数组与链表、栈与队列、顺序表操作、单链表、双向链表、循环队列、括号匹配\\递归与复杂度&时间复杂度、空间复杂度、递归终止条件、递归与分治、循环边界与不变量\\树结构&二叉树、二叉树遍历、二叉搜索树、堆与堆调整、优先队列\\图结构&邻接表与邻接矩阵、图的遍历、深度优先搜索、广度优先搜索、拓扑排序\\排序方法&冒泡排序、快速排序、插入排序、归并排序、排序稳定性\\查找与索引&顺序查找、二分查找、哈希表、哈希冲突处理\\\bottomrule\end{tabularx}\end{table}

先修关系按知识依赖核对。例如，二分查找依赖有序序列上的随机访问和区间缩减思路，并不要求先掌握某一种具体排序算法；括号匹配依赖栈的操作约束；拓扑排序需要图的存储和队列组织，但不把普通图的遍历结果当作唯一前提。

每个知识点配套自编讲义，包含定义、具体例子、操作检查和自测问题。习题覆盖概念、程序步骤、推理与迁移应用，并保留多种题型。数据结构视频展示冒泡排序、栈操作和二分查找，Python视频展示循环边界、列表共享与切片、函数参数与返回值。六段视频的字幕和检查题各自关联相应知识，便于从具体片段继续练习。

内容包以课程代码与稳定样例标识定位实体，正常启动时按版本补齐内容。导入过程记录每个实体与学习事件，重复启动保持已有状态，同时保留教师修订和学生学习历史。这样的组织方式便于逐步补充课程，而不需要重建学生的学习数据。


Python课程按变量与表达式、条件与循环、组合数据、函数、异常与文件、对象与综合问题六个单元展开。主体验学生在两门课程中各有4项目标、12项安排和5段多轮辅导历史；课堂记录包含当前活动、已结束活动和不同学生的回答。学习历史由统一评分函数形成，教师统计直接读取这些记录。

\begin{table}[htbp]\centering\caption{正常初始化后的两门课程内容}\begin{tabularx}{\linewidth}{Yrr}\toprule 内容&数据结构与算法&Python程序设计\\\midrule 章节／知识点／关系边&6／31／38&6／27／30\\已审核题目／题型&39／5&34／6\\讲义／教学方案／视频&31／6／3&29／6／3\\主要学生目标／安排／辅导会话&4／12／5&4／12／5\\课堂活动／学生回答&5／63&4／40\\班级人数&28&14\\\bottomrule\end{tabularx}\end{table}

\subsection{知识图谱的自动布局与交互}

两门演示课程默认呈现按章节聚类的完整知识网络，学生进入后即可看到课程结构与具体知识。章节总览用于快速理解主题构成，局部邻居和先修链用于深入探索。视图偏好按学生与课程分别保存；从特定题目、视频或计划进入时临时聚焦相应知识，普通进入仍保持有效的课程视图。

布局采用ELK分层算法与避障路由\cite{elk,sugiyama}。大图按章节建立复合容器，在章节内部沿先修方向排列，章节之间保留必要关系；局部网络采用较舒展的横向组织。连线先计算避障折线，再在可用通道内处理圆滑转角。圆节点采用一致的连接端口，淡色章节区域帮助辨认归属，选中节点与相关路径共同高亮。
\screen{graph-overview}{六章节总览：先理解课程结构，再展开知识网络}
\screen{graph}{全屏图谱：先修方向、节点状态与详情共同呈现}
节点名称、环形进度和状态文字共同表达学习情况。无评分证据时显示“待学习”，详情显示“待观测”；有证据后再显示估计比例和有效证据量。学生由节点进入视频、练习、资料或辅导，教师由同一对象进入编辑，保存后自动重排并选中新节点。

图谱提供两种不同的阅读操作。“适配全图”根据当前画布与节点边界计算缩放，使完整网络留在视口内；“聚焦阅读”把选中知识放在固定的115%比例，便于读名称和邻近关系。详情可以收起，将空间交还给图谱。全屏隐藏普通导航，保留章节、搜索、关系聚焦、详情和小地图；退出时恢复进入前的相机位置。手机端将详情移到画布下方，支持触摸平移和明确的缩放按钮。

\section{四项核心教学方法}

\subsection{K-TRACE：相关性修正的多维知识证据追踪}
固定加减分会把“独立答对”“多次提示后答对”和“密集重复”视为近似结果。K-TRACE按知识点维护带时间的Beta证据，输入为实际得分、难度、技能、提示、时间与历史作答，输出知识估计、有效证据量、近似区间和复习依据。知识追踪研究说明了按学习过程维护状态的价值\cite{corbett,piech}；本系统选择参数少、计算依据清楚的更新方式，适合课程规模的交互反馈。

设先验为$(\alpha_0,\beta_0)=(1,3)$，距上次有效证据$d$天，半衰期为30天。时间保留比例与折减后参数为
\begin{equation}
r=2^{-d/30},\quad \bar\alpha=1+(\alpha-1)r,\quad
\bar\beta=3+(\beta-3)r .
\end{equation}
设题目难度权重$b\in\{1,1.4,1.8\}$，得分比例$s\in[0,1]$，提示次数$h$，24小时内同知识、同技能的其他题作答数为$m$。独立性质量$q$、相关性修正$c$及权重$w$定义为
\begin{equation}
q=\frac{s}{1+h},\qquad c(m)=\frac{1}{\sqrt{1+0.4m}},\qquad w=b\nu c(m).
\end{equation}
短期重复成功取$\nu=0$，同题其他短期重试取0.35，首次或间隔24小时后的作答取1。相关性修正随相近作答数量增加而减小，避免密集变式被当作完全独立的测量。更新为
\begin{equation}
\alpha'=\bar\alpha+wq,\quad \beta'=\bar\beta+w(1-q),\quad
\hat p=\frac{\alpha'}{\alpha'+\beta'},\quad N'=rN+w .
\end{equation}
\diagram{evidence}{K-TRACE：独立性、相关性与时间共同调节学习证据}
Beta方差为$\alpha'\beta'/[(\alpha'+\beta')^2(\alpha'+\beta'+1)]$，系统使用均值加减1.96倍标准差并截断到0—1作为近似区间\cite{stan}。区间表达信息充分程度；无证据的先验参与后续计算，但不计入已观测知识均值。

概念、步骤、推理与迁移通过题目技能和表现摘要分别保留。先修缺口进入诊断和路径规划，不直接扣减后续知识的作答估计。视频观看与实验操作保存为活动，不转换为独立成功；主观题等待教师确认，65分按$s=0.65$更新，而非将及格统一视为满分证据。

二级独立题权重为1.4；在8次同技能近期其他作答后，权重降为0.68313。同一先验下，独立正确的后验为0.44444，两次提示后正确为0.27160。零权重重复成功保留历史，但不刷新有效证据时间。上述参数为明确的工程设定，课程应用时可结合题目质量逐步标定。

\subsection{PATH-G：先修约束的多目标学习路径规划}
PATH-G把“推荐哪项内容”扩展为“先后完成哪些内容”。输入包含知识有向图、K-TRACE状态、个人目标、教师任务、到期复习、错题和可用资源；输出当次计划、预计用时、后续候选和推荐分项。只有具备已审核可评分题的知识点才形成练习候选，知识卡、视频和实验沿该知识的实际资源组织。

知识$v$已有准备状态的条件设为$N_v\ge2\land \hat p_v\ge0.6$。路径$P=(k_1,\ldots,k_t)$需满足时间预算$B$和先修顺序：
\begin{equation}
\sum_{i=1}^t c_{k_i}\le B,\quad
\forall v\in\mathrm{parents}(k_i),\
v\in\mathrm{ready}\ \lor\ v\in\{k_1,\ldots,k_{i-1}\}.
\end{equation}
前置活动排在前面表示执行顺序，实际作答后再判断是否掌握并重新规划。每项候选的效用由六个明确分量组成：
\begin{equation}
U_k=4(1-\hat p_k)+1.5(u_k-l_k)+3G_k+3T_k+2D_k+\min(3,M_k).
\end{equation}
其中$u-l$为估计区间宽度，$G,T,D$分别表示个人目标、教师任务与到期复习，$M$为未订正错题数。路径目标为
\begin{equation}
J(P)=\sum_{i=1}^{t}\left(0.92^{i-1}U_{k_i}-0.06c_{k_i}\right).
\end{equation}
\diagram{path-search}{PATH-G：筛选可执行候选，在约束内搜索学习顺序}
系统采用束宽24的有限搜索。多个进行中目标共享同一段学习时间，预算取其每日投入的最大值，无进行中目标时取30分钟；最大搜索深度为$\min(|K|,\max(1,B))$，支持一分钟的短活动。每次扩展只接受已满足先修和预算的候选，按目标值、成本和知识编号稳定排序。它比静态排序多了一层序列约束，同时保持本机交互所需的有限计算量。

视频活动用时按实际时长向上取整后增加2分钟思考时间；题目用时按$3+2d+2\mathbf{1}_{\mathrm{多选}}$分钟估计，$d$为难度级别。预算内活动与后续候选分开展示，避免将候选总量误认为今天必须完成的内容。

先修条件同样适用于计划之后的候选排序。某个目标缺少可执行先修资源、且先修尚无充分证据支持时，该目标继续等待，避免出现无法顺序完成的学习建议。学生作答后重新生成路径，先修得到支持便可将当次时间用于后续内容。有限束宽控制搜索规模，活动用时采用明确的估计值；约束满足与动态重规划的对照结果见第六节。

\subsection{DIA-GUIDE：诊断—干预—校验协作教学编排}
普通单轮回答往往止于解释。DIA-GUIDE把解释置于具有输入、职责、工具和反馈的教学过程内。它采用混合代理机制：规则诊断代理读取证据，策略代理选择教学表征并指导一次语言模型生成，规则校验代理核对对象和动作。每轮最多一次模型调用，三个角色通过结构化状态传递结果。

共享状态写为$S=(E,D,P,V,F)$，依次表示学习证据、诊断、策略、校验与理解反馈。诊断输入包括当前题、提示情况、近期错误、先修缺口和上轮检查；输出教学关注和支持证据。策略根据这些结果选择概念对照、渐进提示、实验表征或独立检查，再将课程片段和动作候选传给模型。
\diagram{actions}{DIA-GUIDE：角色职责、一次模型生成与学生反馈}
校验代理只接受当前课程内的动作对象，核对检查题标识、资源归属、支持的动作类型和引用编号。校验结果保存为通过、等待检查、部分可用或需要修正；只有有效理解检查实际存在且校验通过时，才进入等待检查。移除无效动作或引用后保留部分修正状态，学生提交检查后，会话状态和保存记录一起转为已反馈。前端在流式结束时使用已校验文本。通过校验的结果呈现为实验、视频、练习或理解检查入口，学生能够直接执行。

理解检查调用课程已审核客观题，实际评分与普通练习共用证据引擎。会话、角色输出、工具和反馈持久保存；同一知识的前次检查失败时，下一轮优先转换为实验表征，没有相应实验时使用可用的概念对照。流程为
\begin{equation}
E\rightarrow D\rightarrow P\rightarrow V\rightarrow F\rightarrow E_{\mathrm{next}} .
\end{equation}
\screen{tutor}{课程上下文辅导、结构化协作轨迹与理解检查}
自动化案例验证了三个角色的顺序、检查评分、持久反馈以及失败后的策略改变。一次正式DeepSeek抽样也完成协作输出与有效理解检查。相较于纯文本回答，本方法增加了课程内动作执行和反馈状态，帮助辅导从“说完”继续走向“检验并调整”。规则诊断用于可检查的常见困难，开放误解仍由教师结合原始答案判断。教学工具与行动结合的思路亦可见于语言模型行动研究\cite{reactpaper}。

\subsection{CLASS-LOOP：基于错因距离的动态教学分层}
相同总分可能对应不同困难。CLASS-LOOP从最近14天的已评分作答形成六维向量$x_s=(e_c,e_p,e_r,e_t,h,a)$，依次表示概念、步骤、推理和迁移的平均失分、提示比例与参与间隔。缺少某类作答时取中性值0.5；参与间隔$a=\min(1,\Delta d/14)$。输入为班级证据，输出临时群组、代表学生、关注类型、离散度、匹配题目和目标学生。

为避免把“没有观察到”当作“没有困难”，系统为每维增加支持度$c_{sj}$：四类技能取$\min(1,n_{sj}/4)$，提示维度取$\min(1,n_s/6)$，参与维度有评分记录时为1。采用带支持度的加权L1距离和k-medoids聚类\cite{kaufman}：
\begin{equation}
d(s,t)=\sum_{j=1}^{6}w_j[\min(c_{sj},c_{tj})|x_{sj}-x_{tj}|+0.25|c_{sj}-c_{tj}|],\quad
w=(1.2,1.2,1.0,1.1,0.8,0.7).
\end{equation}
共有证据决定错因差异，支持程度不同给予适量惩罚。最近14天少于两次评分作答的学生先进入“学习起点观察组”，通过少量独立练习补充证据。其余学生以确定性的最远点初始化选择代表，每次分配给最近代表，再在组内选取总距离最小的代表；包含观察组在内最多4组，最多更新12轮。目标为
\begin{equation}
\min_{\mathcal C,m_1,\ldots,m_K}
\sum_{k=1}^{K}\sum_{s\in C_k}d(s,m_k).
\end{equation}
\diagram{class-loop}{CLASS-LOOP：错因分组、任务匹配与证据回流}
群组名称取支持度加权中心中最突出的教学关注，未观测维度不参与错因判断。题目匹配先覆盖群组成员的具体薄弱知识，再比较相应技能与难度距离，从当前课程的已审核题中选择3题。相同的步骤困难因而可以分别进入栈操作或Python循环边界的练习。教师可直接带入题目和目标学生，确认后按正常作业流程发布。任务基线保存发布时状态，完成后重新计算群组，避免把临时困难固化为能力标签。

在24名等均分的合成学生中，四类错因原型各6人。方法按主要困难形成四组，原型纯度1.0；按等分编号分组的对照为0.333333。改变单个学生的证据会改变其教学关注，输入顺序重排则保持成员一致。矩阵与原始作答入口帮助教师把分组结果与具体学习经历共同判断。

\subsection{间隔复习与个人学习安排}

复习调度借鉴SuperMemo SM-2的质量等级、难易因子和间隔递增思想\cite{sm2,settles}。平台把独立正确、提示正确和错误分别映射到质量5、3和1。首次成功安排1天后复习，第二次成功安排6天后复习，后续间隔根据难易因子增长。
\begin{align}
E'&=\max\{1.3,E+0.1-(5-Q)[0.08+(5-Q)0.02]\},\\
I_n&=\begin{cases}1,&n=1,\\6,&n=2,\\I_{n-1}E',&n>2.\end{cases}
\end{align}
其中$Q$为本次质量，$E$为难易因子，$I$以天计。错误使连续成功次数归零并回到1天间隔；有提示的成功最多安排2天间隔。短时间的重复成功沿用原复习安排，使反复点击同题不能不断推迟下一次复测。

复习工作区同时呈现错题订正和间隔复习。每项安排显示知识点、下一次时间、间隔和质量依据，学生能够从安排直接进入对应问题。原先答对、但已需要复习的题目也会出现在安排中，避免复习仅围绕错误答案展开。

学生的个人目标和计划提供另一层时间组织。目标关联知识点和每日投入时间。计划保存具体题目、作业、视频资源与片段时间、阅读资料或实验类型，开始学习时直接打开对应对象。已评分作答、有效观看进度与实验步骤分别触发完成判据；手动勾选保存为自报完成。系统推荐与个人安排共同显示，但由实际作答结果决定知识状态。学生因此能够按自己的节奏规划学习，同时保留可解释的复测依据。
\diagram{plan-execution}{学习目标、具体活动与完成事件的连接}
\screen{python-plan}{Python课程的个人计划与预算内活动}


\subsection{连续作答的计算示例}

以一个知识点为例，表中按照同一时刻的独立事件计算更新，专门观察难度、提示和重复作答产生的差异。初始参数为$\alpha=1$、$\beta=3$，知识估计为0.25。题目甲、乙、丙分别代表该知识点下不同题目，最后一行仅推进时间，不新增作答。
\begin{table}[H]\centering\caption{同一知识点的合成证据更新示例}\begin{tabularx}{\linewidth}{Yrrrrr}\toprule 场景&权重&$\alpha$&$\beta$&估计&有效证据\\\midrule
初始先验&0.000&1.000&3.000&0.250&0.000\\
甲：独立答对&1.000&2.000&3.000&0.400&1.000\\
乙：提示后答对&1.400&2.700&3.700&0.422&2.400\\
丙：首次答错&1.400&2.700&5.100&0.346&3.800\\
丙：短期订正&0.490&3.190&5.100&0.385&4.290\\
丙：再次答对&0.000&3.190&5.100&0.385&4.290\\
30天无新证据&0.000&2.095&4.050&0.341&2.145\\
\bottomrule\end{tabularx}\end{table}


首次独立答对题目甲后，正向参数增加1，估计由0.25变为0.40。题目乙的难度权重为1.4，但使用一次提示后答对，其质量为0.5，因此正向与负向参数各增加0.7。系统保留这种差别，使“借助提示完成”和“独立完成”在后续推荐中具有不同含义。

题目丙首次答错会增加负向证据；随后短期订正正确，使用0.35的新颖度系数，记录学生的进步，同时限制同题重复带来的增幅。再次连续答对题目丙时，证据权重为0，估计和有效证据量保持原值。这样，学生能够继续练习，而状态估计主要依靠不同题目与间隔复测形成。

经过30天而没有新证据时，超出先验的两项参数及有效证据量同时折半，估计向先验回落。折减作用于历史证据的分量，因而知识估计通常并不直接减半。表中数值由与平台相同的计算函数生成，属于合成输入下的方法检验。


\section{师生学习与教学应用}

\subsection{学生工作台与课程入口}

学生首页以一项主要学习任务建立页面中心。任务显示知识名称、推荐原因、预计时间和继续入口，视频、资料、计划与课堂入口围绕当前任务组织；趋势、技能与学习历史通过工作视图查看。学生不用先理解平台的全部模块，就可以沿当前需要继续学习。
\screen{student-dashboard}{学生工作台：当前任务优先，课程进度与近期学习相邻}
桌面工作区使用当前视口的剩余空间，主要任务与操作保持可见，大量记录分页，长正文进入阅读详情。页面在浏览器100\%缩放下保留正常字号，课程切换统一放在顶部选择器中。当天没有个人安排时，系统仍依据课程状态提供下一步建议。

课程切换会同步当前课程与对象身份。选择新课程时清理上门课程的题目、资料、会话和图谱引用，取消旧请求，保留当前功能模块；浏览器刷新与前进后退保持课程选择一致。移动端把侧栏和右侧内容调整为适合连续阅读的结构，主要任务仍位于首屏。

学生端采用宋瓷青绿体系。深青瓷用于主要操作和关键路径，纸釉白承载大面积背景，浅青釉和淡青白用于分层。标题以宋体形成平静的阅读节奏，正文和控件保持清晰，减少重复标签和装饰性说明。


知识领域条形图只汇总已有证据的知识；14天趋势区分练习数量和正确率；活动分布保留视频、实验、辅导检查和教师反馈的数量；四技能表现依据已评分题目计算。点击图表后进入实际课程活动，图形与学习入口相邻。

\subsection{练习、反馈与错题复习}

练习界面把题干和作答区域作为视觉中心。选项支持明确的悬浮、选中和禁用状态，渐进提示位于辅助区域。提交后显示确定性评分、解析、错因、知识状态和下一次复习时间。证据反馈区对照作答前后估计，展示提示与重复对本次更新的影响，并提供重新计算的后续学习入口。
\screen{practice}{练习作答结果与证据反馈}
平台支持单选、多选、判断、填空、简答和分步作答。客观题按照预设答案确定性评分；主观作答先保存为待复核状态，教师依据量规确认结果后更新学情。学生能区分即时反馈和等待教师判断的内容。

错题订正保留原来的错误类型、次数和最近答案，订正后仍按质量安排复测。这样，学生能够看到错误如何被纠正，以及需要何时再次独立检验。短时间重做被记为练习经历，但知识状态通过独立性与新颖性规则控制更新。

课程作业把教师安排和个人练习结合起来。学生从任务进入题目，提交记录保留作业标识，教师看到每位学生的完成数量与结果。题库中已经用于作业或已有作答的题目保留答案与知识归属，修改此类内容时通过新建变式维护历史解释的一致性。


\subsection{视频学习与交互实验}

视频工作区以播放器为主体，时间轴提供知识片段，字幕支持搜索和定位，笔记记录学生的思考与时间戳。教师可以编辑字幕、片段和检查点，学生在观看到关键位置时完成随堂检查。
\screen{video}{课程视频、片段导航与时间戳笔记}
六段原创教学视频分别依据算法状态与Python执行状态生成。冒泡排序展示相邻比较和最大值归位；栈操作展示入栈、出栈及空栈边界；二分查找展示中点比较与左右边界更新。每段视频的字幕描述对应实际步骤，检查题关联同一知识点。

交互实验进一步把观察变成主动预测。学生可以修改输入、预测下一步、逐步运行或自动播放，并查看每一步的状态。数据结构实验包含排序、栈、队列、二叉搜索树和广度优先搜索；Python实验使用有限执行预设演示分支、循环、列表与函数调用，逐步展示代码位置和变量变化。实验记录保存到学习事件，供学习轨迹查看；知识状态仍由后续作答补充证据。

\screen{python-experiment}{Python函数实验：代码位置、变量状态与逐步执行}
视频、笔记和辅导之间使用资源标识与时间戳连接。学生在某一段不理解时，可以带着该字幕片段进入辅导，而不需要重新描述整个课程背景。


\subsection{教师课程建设与AI备课}

教师工作台与学生首页采用不同结构。当前课程操作、需要关注的知识、待复核作答和近期学习并列呈现，使教师能够从教学判断直接进入内容维护或任务安排。
\screen{teacher-dashboard}{教师工作台：课程操作、重点知识与待复核作答}
课程管理支持课程、章节、学习目标和班级维护。学生通过班级取得课程访问范围，教师可以使用邀请码或添加已有学生组织教学。章节与知识点分开保存，调整章节时仍可保留相关知识和学习记录。

AI备课以教师输入的知识点和教学意图为起点，生成教学目标、先修知识、教学步骤、例子和理解检查。生成结果进入可编辑草稿，教师修改并确认后保存到课程资源。题目生成同样先经过结构与选项校验，再由教师审核发布。

资源列表采用适合检索和编辑的行式组织。上传资料、修改正文和查看历史版本使用相邻入口，减少重复卡片和说明文案。教师端采用茶烟米褐主题，温润背景与稳定的边界帮助高密度信息保持秩序。


\screen{resources}{教师资料列表：完整标题、自然摘要与稳定操作区}
资料摘要去除Markdown符号后按语义长度呈现，不以固定高度裁切正文；长标题自然换行，查看、编辑和版本操作保持稳定位置。资料按类型筛选、搜索并以每页9项分页，完整正文留在详情与编辑器中。

\subsection{学习证据下钻与针对性教学}

教师学情分为知识矩阵、教学干预、学生轨迹和作答复核四个工作区，切换后保留筛选状态。知识矩阵将28名学生与31个知识点放在一起。每个单元格来自对应学生的K-TRACE状态，显示估计与待观测标记；章节、学生和状态筛选用于缩小问题范围。点击单元格后可查看具体答案、提示、评分方式、权重、质量和前后变化。
\screen{teacher-matrix}{班级知识矩阵：学生差异与原始证据入口}
\screen{teacher-groups}{教学干预工作区：错因群组、匹配题目与任务发布}
趋势图同时呈现作答次数和正确率，状态分布显示同一知识下学生差异，动态群组将困难关联到可发布题目。教师既可以从知识点选择学生，也可以从群组带入名单与题目，从分析直接进入教学安排。

发布任务时保存目标学生的知识状态与有效证据基线。结果页对照完成情况、发布时与当前估计、新增证据和后续得分。主观题单独保留待复核状态，AI按量规提出分项建议，教师确认后只更新一次证据，使任务结果与学习记录相互对应。

\subsection{师生协作的贯通案例}

以“栈与队列”的学习过程为例，学生先完成入栈出栈题，若把入栈顺序误认为出栈顺序，系统记录具体错因。知识图谱中该节点仍关联数组与链表等前置内容，推荐给出先修查看、讲义与针对性练习入口。

学生进入辅导后，当前题目与知识点随提问传入。助手以概念对照解释规则，并提供栈操作实验。学生可以压入1、2、3，再逐次弹出，观察顺序变化。辅导后的理解检查选择同一知识点的另一道已审核题，避免仅复述刚刚看到的原题答案。

检查提交后，实际评分更新证据。若学生独立答对，系统安排后续间隔复测；若仍然错误，下一行动转向概念与先修关系。会话、理解检查和实际作答相互关联，重新打开历史时仍可查看检查结果。

教师从学情页进入该学生的证据，看到原答案、提示次数和状态变化。发布巩固任务时保存基线，学生完成后，教师在结果中对照状态与新证据。一次流程把课程解释、学生练习和教师安排连接起来，体现平台对教学连续性的支持。

\begin{table}[H]\centering\caption{案例中的对象关联}\begin{tabularx}{\linewidth}{p{2.4cm}Y}\toprule 环节&保存与传递的信息\\\midrule 原始练习&学生、知识点、题目、答案、提示与时间\\辅导会话&课程、题目、知识点、资料引用与教学动作\\理解检查&会话、检查题、实际作答及下一行动\\针对性任务&目标学生、题目列表、发布时状态与证据量\\任务结果&完成情况、当前状态与新增证据\\\bottomrule\end{tabularx}\end{table}


\section{数据组织与AI服务运行}

\subsection{课程资料检索与版本化引用}

AI解释使用课程内资料作为事实来源。讲义、字幕和教学资料上传或编辑后，系统保存版本，并通过资源变化队列更新全文检索片段。未变化的资源保留原索引，使重复提问不再导致整门课程资料重新建立索引。
\diagram{retrieval}{资源增量索引与课程范围内检索}
检索采用SQLite FTS5\cite{sqlite,rag}。中文文本按连续二字组合与少量学科单字形成检索词，常见提问词作为停用项处理。问题在当前课程范围内匹配并按相关性排序；指定知识点而检索不足时，使用该知识点关联资料作为补充。

引用包含资源标识、标题、片段、版本和查看入口。学生可以从回答进入资料，教师编辑资料后仍保留历史版本内容。这样的组织使“解释依据是什么”有明确答案，也方便检查教学内容更新后的引用是否同步。

资源触发器统一标记新增、编辑和删除。更新服务只删除变化资源的旧片段，再写入其新片段。测试通过连续检索前后的数据库变化数确认未发生无谓重建，并通过编辑资源后返回的新版本和新片段检查更新结果。


\subsection{数据组织与一致性实现}

课程、章节和知识点构成内容骨架，资源与习题提供学习材料，作答与事件记录学习过程，证据更新与复习安排保存个性化计算的结果。教学任务基线连接任务发布和后续对照。
\diagram{data}{主要数据对象及其关联}
SQLite外键维护对象之间的引用关系，事务把作答保存、状态更新、错题处理和复习安排组织为一次一致操作。主观题提交与教师确认分开保存，界面从复核状态区分两个阶段。

资源版本保存编辑前的正文和元数据，当前资源的版本号随修改递增。题目已经被作业或学习记录引用后，核心答案、题型与知识归属通过新建变式调整，使旧答案和旧学情仍然能够按原来的教学内容解释。

数据库更新使用显式迁移和内容版本。系统升级时保留既有课程、自定义内容和学习历史，自带课程按课程代码、内容版本和稳定实体标识匹配扩充。示例班级通过固定种子生成27名差异化学生，与主体验学生组成28人班级，另设14人的Python班。扩充过程幂等执行，保留已有用户内容与学习历史。


\subsection{运行方式与AI资源控制}

前端使用React 19、TypeScript和Vite 6，后端使用Python 3.13与FastAPI。桌面体验由统一启动入口启动前后端并打开浏览器，重复启动时复用本项目服务，端口冲突时给出明确处理信息。
\diagram{deployment}{本机运行与AI适配层}
常规功能验证采用本地示例模式，使课程、练习、辅导动作和教师任务能够在无需外部接口的条件下复现。正式环境配置后端密钥后，AI适配层自动使用官方服务，接入DeepSeek聊天接口和流式返回\cite{deepseek}，保持相同的会话与资料引用结构。

服务端控制请求频率、输出长度和超时，不自动重试失败请求。真实接口联调通过独立额度记录控制，功能验证与截图均使用本地示例模式。访问凭据只由后端读取，前端页面和公开源码不包含密钥。

前端按路由延迟加载较大的页面\cite{react}。公式与Markdown相关依赖随辅导页面加载，图谱布局和交互实验分别维护。样式按基础控件、学习、教学、响应式、工作区与图谱拆分，双主题通过设计变量覆盖，便于统一调整。


产品交互与开发联调分别记录用量。产品按正常教学请求运行；开发抽样采用单独的四场景账本，无自动重试，已完成后阻止重复执行。正式抽样覆盖辅导、出题、备课和量规评阅，分别验证流式和结构化输出，避免常规回归消耗模型额度。

\section{验证方法与机制对照}

验证分为计算性质、接口业务和浏览器操作三层。计算测试固定时间与输入，检查相关性折减、提示质量、部分得分、时间折减、重复时钟、先修搜索和错因分组。接口测试使用独立数据库，检查作答、主观复核、资源版本、协作教学反馈和多班级汇总。浏览器从师生页面实际发布、完成并查看结果。

后端业务与算法98项、交付校验14项通过；浏览器师生流程与交互回归共23项通过。界面检查覆盖两门课程的师生18个业务模块及登录入口，采用四种桌面视口与手机窄屏，共保存447张截图，页面运行错误为0。自动化、内容建设和录制使用本地示例模式，本次升级新增官方接口调用为0。已有官方服务抽样覆盖辅导、出题、备课与量规评阅，四次均成功、零重试，累计输入3744 Tokens、输出724 Tokens。

\subsection{实验设计与对照条件}

实验采用固定种子4011的隔离合成案例，分别检验更新规律、路径约束、分组区别能力与教学反馈。四项方法使用不同指标，避免用一个总分掩盖各自职责。K-TRACE进行27组消融，PATH-G运行48个图与预算组合，CLASS-LOOP使用24名等均分学生，DIA-GUIDE检查五类策略情景并执行两轮接口对话。合成输入便于重复观察机制差异，实际教学收益将在持续课程使用中通过延迟复测与迁移任务评价。

\begin{table}[htbp]\centering\caption{实验输入、对照与衡量指标}\begin{tabularx}{\linewidth}{p{2cm}YY}\toprule 方法&输入及对照&衡量指标\\\midrule
K-TRACE&3种难度、3种提示、3种相关题数量；移除相应修正&后验估计、更新权重、有效证据量\\
PATH-G&12个8/12/16节点图，4种预算；目标优先静态排序&先修违规步骤、预算违规、重规划顺序\\
CLASS-LOOP&24名等均分学生、4类错因；等分编号分组&原型纯度、平均加权组内距离\\
DIA-GUIDE&5类策略情景；固定渐进提示&预期策略匹配、反馈持久化与动作执行\\\bottomrule\end{tabularx}\end{table}

\subsection{证据修正与路径约束结果}

K-TRACE消融固定其余条件，分别移除提示和相关性修正。所有27组中，提示增加没有抬高独立表现估计，相近题增加没有放大更新权重。二级题在无相近题与8次相近题时，权重分别为1.4和0.68313；同题重复成功权重为0。对已积累12份有效证据的状态，仅推进90天后证据量为1.5、后验0.431818，符合30天半衰期的定义。

PATH-G在48个案例中先修违规0步、预算违规0次，目标优先排序出现36个先修违规步骤。差异来自搜索时逐步检查可执行条件：后续目标的效用再高，也需要先满足基础或在路径中安排前置活动。在16分钟的动态案例中，路径先为$[1,2]$；节点1得到充分作答支持后，路径变为$[2,3]$，将时间释放给后续学习。缺少先修资源的单独案例保持目标阻塞。

\diagram{method-results}{同一合成条件下的路径约束与错因分组对照}

\begin{table}[htbp]\centering\caption{两门课程的个人目标预算检查}\begin{tabular}{lrrr}\toprule 课程&可用分钟&安排分钟&活动数\\\midrule
CS101 & 10 & 10 & 2\\
CS101 & 30 & 30 & 6\\
CS101 & 60 & 60 & 12\\
CS201 & 10 & 10 & 2\\
CS201 & 30 & 30 & 6\\
CS201 & 60 & 58 & 10\\
\bottomrule\end{tabular}\end{table}
在相同课程状态下，只调整进行中目标的每日时间，六种安排全部满足各自预算。短视频与基础练习的成本较低，较短预算仍能形成可执行路径；增加时间后，系统才继续扩展符合先修要求的后续活动。活动完成时，计划记录保留对应事件编号，个人安排与学习经历可以逐项核对。

\subsection{错因分层与协作反馈结果}

CLASS-LOOP的原型纯度为1.0，等分编号分组为0.333333。纯度定义为每组人数最多的错因原型人数之和除以总人数，衡量分组与已知原型的一致程度。两者平均加权组内距离分别为0.01410和1.19528，方法能将等均分背后的不同困难分开。输入顺序重排后成员不变；改变一名学生的证据后，其关注转换为“操作步骤”；近期无评分证据者进入观察组。由此，群组既对证据变化作出响应，也避免将数据缺失直接解释为能力弱项。

DIA-GUIDE在证据不足、概念辨析、先修缺口、独立复测和检查失败五类情景中，预期策略匹配为5/5；固定提示为1/5。两轮本地示例对话先采用概念对照，检查失败后转换到交互实验。接口同时保存三个角色的输出，动作和引用检查通过。该结果说明策略选择、执行与反馈形成了连续过程，匹配率的含义是规则机制覆盖。

28人的隔离班级在本次理解检查后含2195条评分作答，学情快照采用5次SELECT，本次本机测量70.1毫秒。批量读取使知识矩阵、群组与学生轨迹共享一份学习投影，便于教师在七个工作区之间连续操作。

\subsection{图谱与浏览器交互验证}

图谱几何测试覆盖6、31和80节点、四种画布尺寸，节点碰撞和路径穿过无关节点均为0。相机检查包含节点及选中光圈边界，锚点缩放误差小于$10^{-7}$。浏览器逐项检查搜索、章节、全图、115\%阅读视角与全屏恢复。窄屏回归曾定位到旧样式的负边距，清理后全部节点保持在画布内；几何与浏览器检查相互补充。

\section{开发意图控制与可复现性}

\subsection{上下文组织与开发意图控制}
开发首先将课程学习目标拆成可检查的对象与行为，再提供当前模型、接口、已实现功能和修改范围。课程、知识点、资源与题目共用标识；客观评分放在后端；状态更新集中计算；结构变化通过迁移完成。围绕这些约束组织上下文，使界面建设与算法深化沿同一架构推进。

图谱任务提供节点、先修边、课程规模和跨章节案例，验收要求是节点不碰撞、端点位于圆周、搜索与全屏可用；学习引擎任务提供时间、提示、重复与部分得分的反例，要求同一函数输出状态与依据；教学协作任务提供支持的工具和动作集合，要求实际检查结果回到下一轮策略。明确的输入与反例使模型生成的实现可以通过程序检查。
\begin{table}[htbp]\centering\caption{三类开发任务的意图与反馈}\begin{tabularx}{\linewidth}{p{2cm}YY}\toprule 任务&意图控制&实际验证反馈\\\midrule
知识组织&共享课程对象、渐进显示、几何性质&修正章节切换后详情不同步与先验比例显示\\
学习计算&提示、时间、相关性和先修约束&检查部分得分、重复成功时钟和路径预算\\
协作教学&一次生成、结构化消息、业务动作&验证检查失败后改换表征与反馈持久化\\\bottomrule\end{tabularx}\end{table}
生成结果进入代码后，通过单元、接口和浏览器检查反馈。发生失败时记录输入、现象、预期和修改，再进行针对性复测。开发历史按实际功能阶段保留，中文提交说明对应可验证修改；原始指令与实现反馈形成配套协作记录，可与版本演进相互核对。

\subsection{核心实现与快速复现}
知识证据引擎与业务接口分离。以下为状态更新的等价简写，展示公式如何进入实际评分流程；课程与题目筛选在调用前完成。
\begin{lstlisting}[language=Python]
quality = score_ratio / (1 + hints)
correlation = 1 / sqrt(1 + 0.4 * nearby_count)
weight = difficulty_weight * novelty * correlation
alpha = decayed_alpha + weight * quality
beta = decayed_beta + weight * (1 - quality)
value = alpha / (alpha + beta)
\end{lstlisting}
结构化协作先读取诊断状态，再生成课程内解释并校验动作。理解检查将结果关联到会话与实际作答，后续请求读取持久反馈。路径与班级分析批量读取状态、题目和活动，减少同一请求中的重复查询。28人、2195条评分样本的学情投影采用5条SELECT，本次本机测量为70.1毫秒。

本机复现从依赖安装、数据库初始化和统一启动入口进入。没有模型密钥时可通过示例模式检查课程、作业和协作动作；提供有效后端密钥后使用官方服务。推荐的体验顺序为学生图谱与练习、辅导检查、教师矩阵、针对性任务、学生完成和教师结果对照，能够覆盖共享数据的完整教学过程。

\section{设计体系与发展方向}

\subsection{双端设计体系与使用体验}

学脉采用共享排版与控件规范、按角色区分主题的设计方式。学生端以宋瓷青绿营造专注的学习环境，教师端以茶烟米褐组织教学信息。学生深青墨导航与教师深茶褐导航建立角色识别，明亮纸白工作区承载阅读，柔和投影区分主要容器。正文以15px为基础，主标题27px，按钮与表单维持清楚的14至15px。主色集中用于操作、选中和关键数据，减少装饰性标签。
\begin{table}[H]\centering\caption{双端主题与色彩用途}\begin{tabularx}{\linewidth}{p{2cm}p{3cm}Y}\toprule 角色&主色与背景&实际用途\\\midrule 学生&\#2E6A67、\#F1F0E8&主要操作、图谱路径、学习任务与纸白背景\\教师&\#7E6554、\#F5EEE4&课程编辑、任务发布、教学工作区背景\\辅助&\#607E95、\#C6A348、\#6C7F4F&雾蓝区分序列、槐黄强调复习、教师石绿表达积极反馈\\\bottomrule\end{tabularx}\end{table}

\screen{login}{品牌登录入口：书页路径、身份切换与聚焦表单}
登录将原创书页路径图形与表单放入统一容器，外层深色背景建立空间，身份切换自然改变色调。移动端收起大图形，保留品牌、身份与表单，密码显示、键盘焦点与错误反馈均有明确状态。

界面减少不承担操作作用的副标题、口号和重复标签。首页用布局区分当前任务、课程和近期学习；资料采用列表；练习保留题干与作答中心；视频突出播放器和时间线；教师分析用证据与行动相邻的方式组织信息。按钮使用明确动作，保存、加载、权限和服务状态保留必要反馈。

设计检查覆盖1440×1000、1366×768、1920×1080与390×844四种视口。长页面保存全页和重叠分段图，关键表单、结果和异常状态分别检查。颜色之外，节点轮廓、文字、图标和选中边界共同表达状态，便于在不同显示条件下辨认。


\subsection{技术特点与后续发展}

学脉把知识结构、学习证据和教学行动结合在同一过程内。知识图谱不仅提供课程浏览，也决定推荐中的先修路径；作答不仅形成分数，也进入有时间和独立性定义的状态估计；AI不仅解释问题，也通过受支持动作连接实验、视频和理解检查；教师发布的任务则保留基线，使后续表现能够回到原来的教学判断中。

平台的可解释性来自具体数据，而非额外生成一段理由。一次状态更新有对应题目、时间、提示和权重，一次推荐有先修和可执行步骤，一次教学干预有发布时状态和后续结果。这些信息使学生和教师能够理解建议的来源，并根据实际教学经验调整参数和内容。

当前实现以课程和班级为组织单位，可随教学内容扩展应用场景。后续工作将重点开展题目质量分析与模型参数标定，观察不同题目对同一能力的区分程度；完善延迟复测与迁移任务，扩展知识状态随时间变化的验证；在经过组织和授权的课程实践中，设计更长周期的教学评价，比较学习行为、复测表现与教师使用体验。

内容方面，将继续扩充完整课程案例、较长教学视频和面向不同难度的题目变式。运行方面，将完善正式账号生命周期、数据备份与部署维护，使教学平台能够从本机作品体验逐步走向持续课程使用。


\FloatBarrier

\begin{thebibliography}{99}\setstretch{1.12}\setlength{\itemsep}{4pt}\addcontentsline{toc}{section}{参考文献}
\bibitem{bloom}Bloom B S. The 2 Sigma Problem: The Search for Methods of Group Instruction as Effective as One-to-One Tutoring[J]. Educational Researcher, 1984, 13(6): 4--16. DOI:10.3102/0013189X013006004.
\bibitem{hattie}Hattie J, Timperley H. The Power of Feedback[J]. Review of Educational Research, 2007, 77(1): 81--112. DOI:10.3102/003465430298487.
\bibitem{vanlehn}VanLehn K. The Relative Effectiveness of Human Tutoring, Intelligent Tutoring Systems, and Other Tutoring Systems[J]. Educational Psychologist, 2011, 46(4): 197--221. DOI:10.1080/00461520.2011.611369.
\bibitem{corbett}Corbett A T, Anderson J R. Knowledge tracing: Modeling the acquisition of procedural knowledge[J]. User Modeling and User-Adapted Interaction, 1994, 4: 253--278. DOI:10.1007/BF01099821.
\bibitem{piech}Piech C, Bassen J, Huang J, et al. Deep Knowledge Tracing[C]//Advances in Neural Information Processing Systems 28. 2015: 505--513. \url{https://arxiv.org/abs/1506.05908}.
\bibitem{settles}Settles B, Meeder B. A Trainable Spaced Repetition Model for Language Learning[C]//Proceedings of ACL. 2016: 1848--1858. DOI:10.18653/v1/P16-1174.
\bibitem{sugiyama}Sugiyama K, Tagawa S, Toda M. Methods for Visual Understanding of Hierarchical System Structures[J]. IEEE Transactions on Systems, Man, and Cybernetics, 1981, 11(2): 109--125. DOI:10.1109/TSMC.1981.4308636.
\bibitem{kaufman}Kaufman L, Rousseeuw P J. Finding Groups in Data: An Introduction to Cluster Analysis[M]. New York: Wiley, 1990. DOI:10.1002/9780470316801.
\bibitem{rag}Lewis P, Perez E, Piktus A, et al. Retrieval-Augmented Generation for Knowledge-Intensive NLP Tasks[C]//Advances in Neural Information Processing Systems 33. 2020: 9459--9474. \url{https://arxiv.org/abs/2005.11401}.
\bibitem{reactpaper}Yao S, Zhao J, Yu D, et al. ReAct: Synergizing Reasoning and Acting in Language Models[C]//ICLR. 2023. \url{https://arxiv.org/abs/2210.03629}.
\bibitem{elk}Eclipse Foundation. Eclipse Layout Kernel: ELK Layered[EB/OL]. [2026-11-11]. \url{https://eclipse.dev/elk/reference/algorithms/org-eclipse-elk-layered.html}.
\bibitem{stan}Stan Development Team. Continuous Distributions on [0,1]: Beta Distribution[EB/OL]. [2026-11-11]. \url{https://mc-stan.org/docs/functions-reference/continuous_distributions_on_0_1.html}.
\bibitem{sm2}Wozniak P. SuperMemo 2: Algorithm[EB/OL]. [2026-11-11]. \url{https://super-memory.com/english/ol/sm2.htm}.
\bibitem{sqlite}SQLite Development Team. SQLite FTS5 Extension[EB/OL]. [2026-11-11]. \url{https://www.sqlite.org/fts5.html}.
\bibitem{deepseek}DeepSeek. Chat Completions API[EB/OL]. [2026-11-11]. \url{https://api-docs.deepseek.com/api/create-chat-completion}.
\bibitem{react}Meta Open Source. lazy: React Reference[EB/OL]. [2026-11-11]. \url{https://react.dev/reference/react/lazy}.
\end{thebibliography}
\end{document}
